\section{Overall Framework}

LectureMind system process is divided into multiple parts. Firstly, an automatically extracted text feature is generated by the method of sound acquisition; then, a combination of image features obtained via frame level analysis. Neither signal is sufficient by itself. The transcript refers to what the teacher says, and the slide streaming serves as a generalisation for topic variations. Combine them to generate chunks more aligned with the transition of semantics in lectures. The above three items then form the building blocks of knowledge acquisition in the following steps.

High-level division of the pipeline at this stage includes four steps.

\begin{itemize}
\item \textbf{Multimodal preprocessing}: derive transcript and visual transition signals from the original video;
\item \textbf{Semantic chunking}: combine visual and speech cues to partition the lecture into coherent sections;
\item \textbf{Knowledge organization}: extract fine-grained knowledge points and aggregate them into higher-level structures;
\item \textbf{Grounded interaction}: answer student questions using direct, retrieval-based, or agentic strategies.
\end{itemize}

To make the workflow easier to implement, the pipeline is organized as a dependency-aware task graph. Figure~\ref{fig:task-dag} shows the original 10-task DAG used in LectureMind for preprocessing, knowledge extraction, and summarization.

\begin{figure}[H]
    \centering
    \begin{tikzpicture}[
        node distance=1.3cm and 1.1cm,
        task/.style={rectangle, draw, rounded corners, minimum width=2.45cm, minimum height=0.72cm, align=center, font=\footnotesize},
        arrow/.style={->, >=stealth, thick}
    ]
        \node[task, fill=blue!10] (t2) at (0,0) {T2: HLS\\Encoding};
        \node[task, fill=blue!10, left=of t2] (t1) {T1: ASR\\Transcription};
        \node[task, fill=blue!10, right=of t2] (t3) {T3: SSIM\\Detection};

        \node[task, fill=green!10, below=of t2] (t4) {T4: Thumbnail\\Generation};
        \node[task, fill=teal!10, below of=t4] (t5) {T5: Slide\\OCR};
        \node[task, fill=yellow!10, below of=t5] (t6) {T6: Hybrid\\Chunking};
        \node[task, fill=orange!10, below of=t6] (t7) {T7: Fine-Grained\\Knowledge};
        \node[task, fill=red!10, below of=t7] (t8) {T8: Embed\\Knowledge};
        \node[task, fill=purple!10, below of=t8] (t9) {T9: Coarse\\Summary};
        \node[task, fill=magenta!10, below of=t9] (t10) {T10: Mindmap};

        \draw[arrow] (t3.south) |- (t4.east);
        \draw[arrow] (t2) -- (t4);
        \draw[arrow] (t1.south) |- (t4.west);
        \draw[arrow] (t4) -- (t5);
        \draw[arrow] (t5) -- (t6);
        \draw[arrow] (t6) -- (t7);
        \draw[arrow] (t7) -- (t8);
        \draw[arrow] (t8) -- (t9);
        \draw[arrow] (t9) -- (t10);

        \node[left=1.2cm of t4, font=\scriptsize, align=center] {T1, T2, T3\\run in parallel};
        \node[right=0.9cm of t6, font=\scriptsize, align=left] {T6 reads T1\\output from DB};
    \end{tikzpicture}
    \caption{Task Dependency Graph (DAG) --- 10 Tasks}
    \label{fig:task-dag}
\end{figure}

Three design choices run through the whole framework. Slide transition and speech time-out are considered together but not alternatives; The lecturer's lecture plan cannot be completely obtained from any of these modalities. Secondly, hierarchical knowledge Representation is used. Raw transcription sentences are usually not long enough for teachers to fully grasp the lecture content; whereas an all-encompassing generalisation fails precisely with respect to a particular user's question. The third is that the response strategy adapts flexibly; According to different questions, it may respond directly immediately retrieve data quickly switch to an adaptive-agential mode. Mostly an actual judgment; generally speaking, queries do not require all-optimal paths at once.